\section{Mettre un système bouclé sous forme d’un schéma-bloc}
\begin{enumerate}
    \item Déterminer la fonction de transfert de chacune des parties du système.
    \item Représenter un schéma bloc en faisant apparaitre ces fonctions de transfert.
\end{enumerate}

\section{Déterminer la fréquence d’oscillation d’un oscillateur sinusoïdal}
\begin{enumerate}
    \item A partir du schéma bloc ou des fonctions de transfert, écrire un système de deux équations reliant les deux tensions du système bouclé.
    \item Combiner les équations en remplaçant une des tensions par son expression et simplifier pour ne plus avoir de tension.
    \item Passer en complexe si on est dans le domaine de Laplace.
    \item Prendre partie réelle et partie imaginaire.
\end{enumerate}

\section{Déterminer la condition d’oscillation sinusoïdales}
\begin{enumerate}
    \item A partir du schéma bloc ou des fonctions de transfert, écrire un système de deux équations reliant les deux tensions du système bouclé.
    \item Combiner les équations en remplaçant une des tensions par son expression et simplifier pour ne plus avoir de tension.
    \item Passer en complexe si on est dans le domaine de Laplace.
    \item Prendre partie réelle et partie imaginaire.
\end{enumerate}

\section{Déterminer les conditions de démarrage des oscillations}
\begin{enumerate}
    \item A partir du schéma bloc ou des fonctions de transfert, écrire un système de deux équations reliant les deux tensions du système bouclé.
    \item Combiner les équations en remplaçant une des tensions par son expression.
    \item Passer dans le domaine temporel pour obtenir une équation différentielle sur une tension.
    \item Écrire la condition de stabilité (le système ne doit \textbf{pas} être stable.)
\end{enumerate}

\section{Déterminer l’amplitude des différents signaux mesurables pour un oscillateur quasi-sinusoïdal}
\begin{enumerate}
    \item Identifier les tensions qui sont des sorties d'ALI. Ces tensions sont bornées par $V_{sat}$.
    \item Pour les autres tensions, utiliser la fonction de transfert des blocs dont elles sont la sortie.
\end{enumerate}

\section{Déterminer la forme des signaux et leur fréquence pour un oscillateur à relaxation}
\begin{enumerate}
    \item Identifier les tensions sortant d'un ALI en régime instable. Ces tensions sont des créneaux.
    \item A partir des équations différentielles régissant les autres blocs, établir les formes des autres signaux.
\end{enumerate}